\documentclass[10pt]{article}
\usepackage[a4paper,margin=2.1cm]{geometry}
\usepackage{helvet}
\renewcommand{\familydefault}{\sfdefault}
\usepackage{parskip}
\usepackage{enumitem}
\usepackage[hidelinks]{hyperref}
\setlength{\parskip}{5pt}
\setlist[itemize]{topsep=2pt,itemsep=2pt,parsep=0pt,leftmargin=1.4em}
\pagestyle{empty}

\begin{document}

\noindent 2 July 2026

\vspace{4pt}
\noindent To the Editorial Office\\
\emph{Materials} (MDPI)

\vspace{10pt}
\noindent \textbf{Re: Submission of an original research manuscript}

\vspace{6pt}
\noindent Dear Editors,

We are pleased to submit our manuscript entitled \textbf{``Interpretable Machine
Learning for One-Part Fly-Ash/Slag Geopolymer Strength Prediction: Toward
Multifunctional Binder Design''} for consideration as an Article in \emph{Materials}.

One-part (``just-add-water'') geopolymers replace the hazardous liquid activators of
conventional systems with a dry, pre-blended solid activator, reproducing the ordinary
Portland cement use-model while cutting embodied carbon substantially. This makes them
uniquely suited to field-deployable, low-carbon construction---including protective and
civil-infrastructure applications---but their high-dimensional formulation space defeats
the one-factor-at-a-time experimentation that still dominates the literature. Our
manuscript reviews the materials science of one-part geopolymers, contributes a new
comparative and interpretable machine-learning (ML) analysis, and proposes a
generalizable AI-assisted design framework.

The principal contributions that we believe make this work a good fit for \emph{Materials}
are:

\begin{itemize}
\item A comparative evaluation of four regression models on a published 80-mixture
one-part fly-ash/GGBS dataset under \emph{both} a conventional random split and a
stricter \emph{leave-one-source-out} (LOSO) cross-validation---the honest test of transfer
to a new laboratory---with LOSO reported as the headline result.
\item An explicit treatment of precursor collinearity: because fly-ash and GGBS contents
are near-perfectly anti-correlated ($r=-0.99$), we collapse the precursor axis to a single
slag-fraction descriptor, removing the instability that would otherwise corrupt feature
attribution.
\item SHAP-based interpretation that recovers mechanistically coherent strength
levers---precursor balance and activator Na\textsubscript{2}O dosage---directly from held-out
data, yielding actionable design guidance rather than black-box predictions.
\item A five-stage, closed-loop design framework that clearly delineates what is
\emph{demonstrated} in this study from what is \emph{proposed} as future work, avoiding the
over-claiming common in this area.
\item Full reproducibility: all analysis code, the cleaned dataset with per-mix
source-study identifiers, and the bibliometric corpus and screening tables are archived on
Zenodo (DOI \href{https://doi.org/10.5281/zenodo.21138235}{10.5281/zenodo.21138235}).
\end{itemize}

The analyzed dataset originates from an open-access (CC~BY) publication and is fully
attributed in the manuscript; the modelling, LOSO evaluation, collinearity treatment, and
interpretation are our own original contributions.

We confirm that this manuscript is original, has not been published previously, and is not
under consideration for publication elsewhere. All authors have read and approved the
submitted version and agree to its submission to \emph{Materials}. The authors declare no
conflicts of interest.

\vspace{4pt}
\noindent Correspondence regarding this submission should be directed to the corresponding
author.

\vspace{14pt}
\noindent Sincerely,

\vspace{4pt}
\noindent \textbf{Vinoth Nageshwaran} (corresponding author)\\
School of Computer and Information Sciences, University of the Cumberlands,
Williamsburg, KY, USA\\
\href{mailto:vnageshwaran41452@ucumberlands.edu}{vnageshwaran41452@ucumberlands.edu}

\vspace{6pt}
\noindent on behalf of the co-authors Sudhir Amritphale (Alchemy Geopolymer Solutions,
Ruston, LA, USA) and Soundararajan Ezekiel (Indiana University of Pennsylvania,
Indiana, PA, USA).

\end{document}
